% Options for packages loaded elsewhere
\PassOptionsToPackage{unicode}{hyperref}
\PassOptionsToPackage{hyphens}{url}
%
\documentclass[
]{article}
\usepackage{amsmath,amssymb}
\usepackage{lmodern}
\usepackage{iftex}
\ifPDFTeX
  \usepackage[T1]{fontenc}
  \usepackage[utf8]{inputenc}
  \usepackage{textcomp} % provide euro and other symbols
\else % if luatex or xetex
  \usepackage{unicode-math}
  \defaultfontfeatures{Scale=MatchLowercase}
  \defaultfontfeatures[\rmfamily]{Ligatures=TeX,Scale=1}
\fi
% Use upquote if available, for straight quotes in verbatim environments
\IfFileExists{upquote.sty}{\usepackage{upquote}}{}
\IfFileExists{microtype.sty}{% use microtype if available
  \usepackage[]{microtype}
  \UseMicrotypeSet[protrusion]{basicmath} % disable protrusion for tt fonts
}{}
\makeatletter
\@ifundefined{KOMAClassName}{% if non-KOMA class
  \IfFileExists{parskip.sty}{%
    \usepackage{parskip}
  }{% else
    \setlength{\parindent}{0pt}
    \setlength{\parskip}{6pt plus 2pt minus 1pt}}
}{% if KOMA class
  \KOMAoptions{parskip=half}}
\makeatother
\usepackage{xcolor}
\IfFileExists{xurl.sty}{\usepackage{xurl}}{} % add URL line breaks if available
\IfFileExists{bookmark.sty}{\usepackage{bookmark}}{\usepackage{hyperref}}
\hypersetup{
  pdftitle={Medición de la Interdisciplinariedad: Un Panel de Indicadores Multicomponente para la Evaluación de la Investigación},
  pdfauthor={A. Rivero},
  hidelinks,
  pdfcreator={LaTeX via pandoc}}
\urlstyle{same} % disable monospaced font for URLs
\usepackage{longtable,booktabs,array}
\usepackage{calc} % for calculating minipage widths
% Correct order of tables after \paragraph or \subparagraph
\usepackage{etoolbox}
\makeatletter
\patchcmd\longtable{\par}{\if@noskipsec\mbox{}\fi\par}{}{}
\makeatother
% Allow footnotes in longtable head/foot
\IfFileExists{footnotehyper.sty}{\usepackage{footnotehyper}}{\usepackage{footnote}}
\makesavenoteenv{longtable}
\setlength{\emergencystretch}{3em} % prevent overfull lines
\providecommand{\tightlist}{%
  \setlength{\itemsep}{0pt}\setlength{\parskip}{0pt}}
\setcounter{secnumdepth}{-\maxdimen} % remove section numbering
\ifLuaTeX
  \usepackage{selnolig}  % disable illegal ligatures
\fi

\title{Medición de la Interdisciplinariedad: Un Panel de Indicadores
Multicomponente para la Evaluación de la Investigación}
\author{A. Rivero}
\date{2026}

\begin{document}
\maketitle
\begin{abstract}
Revisamos el panorama de indicadores bibliométricos para medir
interdisciplinariedad y proponemos un panel estructurado de tres
componentes que combina diversidad Rao-Stirling, coherencia de red
(fuerza media de enlace) y una aproximación de efecto transcampo. A
partir de resultados recientes sobre baja consistencia y validez
limitada de enfoques monoindicador, organizamos la literatura en una
taxonomía de cuatro dimensiones conceptuales --- diversidad, coherencia,
difusión y novedad --- y cuatro familias metodológicas. Con un conjunto
de datos universitario de juguete, mostramos que ningún escalar único
distingue integración interdisciplinaria genuina de amplitud polímata o
especialización estrecha. El panel completo sí lo hace y mantiene esa
capacidad bajo perturbaciones de la matriz de similitud. Un estudio de
caso departamental confirma la utilidad práctica del enfoque para
decisiones institucionales y de agencia nacional.
\end{abstract}

\hypertarget{introducciuxf3n}{%
\section{Introducción}\label{introducciuxf3n}}

La investigación interdisciplinaria (IDR) se considera estratégica para
problemas científicos y sociales complejos. Sin embargo, medirla sigue
siendo difícil. La literatura bibliométrica acumula décadas de
propuestas, pero no convergió en un indicador único robusto.

Wang y Schneider (2020), al contrastar 23 medidas en cuatro familias
metodológicas, encontraron baja correlación incluso entre medidas
diseñadas para capturar dimensiones similares. Leydesdorff, Wagner y
Bornmann (2019) mostraron que el índice Rao-Stirling puede verse
dominado por su componente de disparidad, con baja potencia
discriminante en algunos contextos. Cantone (2024) sostiene que la
interdisciplinariedad es un constructo polisémico imposible de comprimir
fielmente en un solo número.

Este trabajo aporta dos contribuciones. Primero, una revisión
estructurada del área mediante una taxonomía de cuatro dimensiones
(diversidad, coherencia, difusión, novedad) y cuatro familias
metodológicas (referencias, citación, semántica textual, red). Segundo,
un panel operativo de tres componentes --- diversidad, coherencia y
efecto transcampo --- con análisis de discriminación, robustez y
aplicabilidad evaluativa.

\hypertarget{fundamentos-conceptuales}{%
\section{Fundamentos Conceptuales}\label{fundamentos-conceptuales}}

\hypertarget{marco-de-diversidad-de-stirling}{%
\subsection{Marco de Diversidad de
Stirling}\label{marco-de-diversidad-de-stirling}}

El problema de medir interdisciplinariedad es, en su núcleo, un problema
de medición de diversidad. Stirling (2007) distingue tres propiedades:
\emph{variedad}, \emph{balance} y \emph{disparidad}. Índices clásicos
como Shannon o Herfindahl capturan sobre todo variedad y balance, pero
no distancia cognitiva entre categorías.

La formulación Rao-Stirling integra explícitamente esas tres
dimensiones:

\[\Delta = \sum_{i \neq j} d_{ij} p_i p_j, \qquad d_{ij}=1-s_{ij}.\]

Aquí \(p_i\) es la proporción en categoría \(i\) y \(s_{ij}\) la
similitud entre categorías \(i,j\).

\hypertarget{evoluciuxf3n-histuxf3rica}{%
\subsection{Evolución Histórica}\label{evoluciuxf3n-histuxf3rica}}

La formalización moderna de IDR acelera tras la tipología OECD (1998) y
los trabajos de Morillo, Bordons y Gómez (2001, 2003). Esa tradición
separa multidisciplinariedad (yuxtaposición), interdisciplinariedad
(integración) y transdisciplinariedad (marcos que trascienden
fronteras).

Wagner et al.~(2011) subrayan además que en la práctica suele
confundirse la dirección de entrada (insumos disciplinares del diseño de
investigación) con la de salida (distribución disciplinar de
publicaciones e impacto).

\hypertarget{problema-de-polisemia}{%
\subsection{Problema de Polisemia}\label{problema-de-polisemia}}

Interdisciplinariedad no es un concepto único. Puede describir
diversidad de insumos cognitivos, estructura social de colaboración o
hibridación semántica del contenido textual. Estas dimensiones no
necesariamente se mueven juntas, lo que explica discordancias entre
autoidentificación y métricas bibliométricas (Zwanenburg, 2022; Aksnes,
Karlstrøm y Piro, 2026).

\hypertarget{desafuxedos-epistemoluxf3gicos}{%
\subsection{Desafíos
Epistemológicos}\label{desafuxedos-epistemoluxf3gicos}}

Los indicadores no solo miden, también performan: redefinen qué se
considera ``buena'' investigación en contextos evaluativos (Rafols,
2019). Por tanto, el diseño de indicadores debe explicitar supuestos y
límites de validez, no solo optimizar facilidad computacional.

\hypertarget{validez-conceptual-vs-empuxedrica}{%
\subsection{Validez Conceptual vs
Empírica}\label{validez-conceptual-vs-empuxedrica}}

Un indicador puede ser matemáticamente elegante y, aun así, mal anclado
al constructo de interés. En IDR esto ocurre cuando una métrica premia
amplitud sin integración o mezcla señales conceptualmente distintas en
un único valor.

\hypertarget{pluralismo-metodoluxf3gico}{%
\subsection{Pluralismo Metodológico}\label{pluralismo-metodoluxf3gico}}

El campo requiere pluralismo disciplinado: diferentes herramientas para
preguntas distintas, con fronteras interpretativas claras y protocolos
de calibración explícitos.

\hypertarget{taxonomuxeda-de-indicadores-de-interdisciplinariedad}{%
\section{Taxonomía de Indicadores de
Interdisciplinariedad}\label{taxonomuxeda-de-indicadores-de-interdisciplinariedad}}

\hypertarget{indicadores-de-diversidad}{%
\subsection{Indicadores de Diversidad}\label{indicadores-de-diversidad}}

\hypertarget{marco-triuxe1dico-de-stirling}{%
\subsubsection{Marco triádico de
Stirling}\label{marco-triuxe1dico-de-stirling}}

La tríada variedad-balance-disparidad organiza la mayor parte de la
bibliometría de IDR.

\hypertarget{medidas-de-variedad-y-balance}{%
\subsubsection{Medidas de variedad y
balance}\label{medidas-de-variedad-y-balance}}

Shannon: \[H = -\sum_i p_i \log p_i\]

Simpson/Herfindahl complementario: \[D_{\mathrm{Sim}} = 1-\sum_i p_i^2\]

Son útiles, pero no incorporan distancia cognitiva explícita.

\hypertarget{uxedndice-rao-stirling}{%
\subsubsection{Índice Rao-Stirling}\label{uxedndice-rao-stirling}}

\[\Delta = \sum_{i \neq j} d_{ij} p_i p_j\]

Integra variedad, balance y disparidad, pero su comportamiento depende
de la matriz de similitud, la taxonomía y la transformación
distancia/similitud.

\hypertarget{diversidad-verdadera-tipo-hill}{%
\subsubsection{Diversidad verdadera tipo
Hill}\label{diversidad-verdadera-tipo-hill}}

Hill (1973), Jost (2006, 2009) y Leinster-Cobbold (2012) enfatizan la
escala de números efectivos:

\[^qD = \left(\sum_i p_i^q\right)^{1/(1-q)}\]

Esta familia facilita comparaciones interpretables y satisface
propiedades de replicación que no siempre cumplen las entropías crudas.

\hypertarget{multi-asignaciuxf3n-y-tipologuxedas-tempranas}{%
\subsubsection{Multi-asignación y tipologías
tempranas}\label{multi-asignaciuxf3n-y-tipologuxedas-tempranas}}

Morillo et al.~(2001, 2003) mostraron que la multi-asignación de
revistas captura señales útiles de hibridación de áreas, aunque con
resolución limitada para distinguir integración profunda.

\hypertarget{estimaciuxf3n-de-matriz-de-similitud}{%
\subsubsection{Estimación de matriz de
similitud}\label{estimaciuxf3n-de-matriz-de-similitud}}

La construcción clásica usa flujos de citación y cosenos (Wang y
Schneider, 2020), en línea con tradiciones de mapeo JCR (Leydesdorff,
2006). Decisiones como Salton vs Ochiai o \(d_{ij}=1-s_{ij}\) vs
\(d_{ij}=1/s_{ij}\) (Jensen y Lutkouskaya, 2014) cambian fuertemente
resultados.

\hypertarget{normalizaciuxf3n-calibraciuxf3n-y-granularidad}{%
\subsubsection{Normalización, calibración y
granularidad}\label{normalizaciuxf3n-calibraciuxf3n-y-granularidad}}

Una misma cartera puede cambiar de clasificación al pasar de esquemas
gruesos a finos (por ejemplo, 40 vs 250 categorías). Esto obliga a
reportar sensibilidad y no tratar valores puntuales como invariantes.

\hypertarget{evaluaciuxf3n-de-validez-marco-zwanenburg}{%
\subsubsection{Evaluación de validez (marco
Zwanenburg)}\label{evaluaciuxf3n-de-validez-marco-zwanenburg}}

Zwanenburg, Nakhoda y Whigham (2022) proponen criterios explícitos de
consistencia y validez. Su lectura principal: no existe una medida única
que satisfaga simultáneamente todos los criterios relevantes para IDR.

\hypertarget{diversidad-de-autor-vs-diversidad-de-referencias}{%
\subsubsection{Diversidad de autor vs diversidad de
referencias}\label{diversidad-de-autor-vs-diversidad-de-referencias}}

Abramo, D'Angelo y Zhang (2018) muestran que la diversidad del equipo
autoral y la diversidad de base referencial capturan capas distintas del
fenómeno.

\hypertarget{resumen}{%
\subsubsection{Resumen}\label{resumen}}

Diversidad es indispensable, pero por sí sola no distingue integración
interdisciplinaria de amplitud polímata.

\hypertarget{indicadores-de-coherencia}{%
\subsection{Indicadores de Coherencia}\label{indicadores-de-coherencia}}

\hypertarget{operacionalizaciuxf3n-por-acoplamiento-bibliogruxe1fico}{%
\subsubsection{Operacionalización por acoplamiento
bibliográfico}\label{operacionalizaciuxf3n-por-acoplamiento-bibliogruxe1fico}}

Rafols y Meyer (2009) introducen coherencia como conectividad interna de
la cartera:

\[S = \frac{1}{\binom{n}{2}}\sum_{k<l}\cos(\mathbf r_k,\mathbf r_l).\]

\hypertarget{evidencia-de-ortogonalidad-respecto-a-diversidad}{%
\subsubsection{Evidencia de ortogonalidad respecto a
diversidad}\label{evidencia-de-ortogonalidad-respecto-a-diversidad}}

La literatura encuentra que diversidad y coherencia no son redundantes.
Perfiles con \(\Delta\) alta pueden tener \(S\) muy baja (polímatas),
mientras integradores muestran simultáneamente amplitud y conectividad.

\hypertarget{marco-diversidad-coherencia}{%
\subsubsection{Marco
diversidad-coherencia}\label{marco-diversidad-coherencia}}

La lectura conjunta (\(\Delta\), \(S\)) permite separar amplitud de
integración, mejorando interpretabilidad frente a escalares únicos.

\hypertarget{alternativas}{%
\subsubsection{Alternativas}\label{alternativas}}

Se han explorado centralidades (betweenness), clustering y otros rasgos
de red; suelen requerir controles de tamaño y normalización para no
confundir posición estructural con interdisciplinariedad sustantiva.

\hypertarget{consideraciones-computacionales}{%
\subsubsection{Consideraciones
computacionales}\label{consideraciones-computacionales}}

Carteras grandes incrementan costo cuadrático en comparaciones
par-a-par; la implementación debe contemplar eficiencia y
reproducibilidad.

\hypertarget{resumen-1}{%
\subsubsection{Resumen}\label{resumen-1}}

Coherencia aporta la dimensión integrativa ausente en gran parte de
métricas de diversidad.

\hypertarget{indicadores-de-difusiuxf3n}{%
\subsection{Indicadores de Difusión}\label{indicadores-de-difusiuxf3n}}

\hypertarget{fundamentos-conceptuales-1}{%
\subsubsection{Fundamentos
conceptuales}\label{fundamentos-conceptuales-1}}

Difusión evalúa si la producción impacta fuera del campo principal. Es
una propiedad de salida y no reemplaza diversidad de insumos.

\hypertarget{definiciuxf3n-operacional}{%
\subsubsection{Definición operacional}\label{definiciuxf3n-operacional}}

Usamos:

\[E = \frac{\text{citas recibidas fuera de la categoría principal}}{\text{citas totales recibidas}}.\]

\hypertarget{dinuxe1mica-temporal}{%
\subsubsection{Dinámica temporal}\label{dinuxe1mica-temporal}}

Ventanas cortas de citación sesgan resultados; la comparación entre
etapas de carrera exige normalización temporal.

\hypertarget{patrones-empuxedricos}{%
\subsubsection{Patrones empíricos}\label{patrones-empuxedricos}}

Resultados recientes (incluyendo Xiang, Romero y Teplitskiy, 2025)
muestran que amplitud temática e impacto transcampo no siempre se mueven
juntos.

\hypertarget{distinciuxf3n-con-interdisciplinariedad-temuxe1tica}{%
\subsubsection{Distinción con interdisciplinariedad
temática}\label{distinciuxf3n-con-interdisciplinariedad-temuxe1tica}}

Una cartera puede ser temáticamente amplia pero con difusión limitada
fuera de sus comunidades de publicación.

\hypertarget{relaciuxf3n-con-integraciuxf3n}{%
\subsubsection{Relación con
integración}\label{relaciuxf3n-con-integraciuxf3n}}

Difusión complementa diversidad y coherencia; no las sustituye.

\hypertarget{validez-y-cautelas}{%
\subsubsection{Validez y cautelas}\label{validez-y-cautelas}}

Diferencias de cultura de citación entre campos obligan a normalizar por
área.

\hypertarget{guuxeda-pruxe1ctica}{%
\subsubsection{Guía práctica}\label{guuxeda-pruxe1ctica}}

Reportar difusión junto con incertidumbre y contexto disciplinar.

\hypertarget{indicadores-de-novedad}{%
\subsection{Indicadores de Novedad}\label{indicadores-de-novedad}}

\hypertarget{fundamento}{%
\subsubsection{Fundamento}\label{fundamento}}

Novedad captura combinaciones atípicas o aparición temprana de enlaces
entre campos.

\hypertarget{enfoques-por-divergencia-estaduxedstica}{%
\subsubsection{Enfoques por divergencia
estadística}\label{enfoques-por-divergencia-estaduxedstica}}

Familias tipo Uzzi (2013) usan comparación observado-vs-esperado en
pares de revistas/referencias.

\hypertarget{muxe9todo-de-permutaciuxf3n}{%
\subsubsection{Método de
permutación}\label{muxe9todo-de-permutaciuxf3n}}

Los modelos nulos son sensibles a supuestos estructurales y a cobertura
de red.

\hypertarget{novedad-temporal}{%
\subsubsection{Novedad temporal}\label{novedad-temporal}}

Wang et al.~(2017) premian combinaciones tempranas, pero su
interpretación puede confundir innovación con rareza transitoria.

\hypertarget{dependencia-de-la-referencia-base}{%
\subsubsection{Dependencia de la referencia
base}\label{dependencia-de-la-referencia-base}}

La elección de referencia esperada (uniforme, permutada, jerárquica)
cambia el significado de la métrica.

\hypertarget{exclusiuxf3n-del-panel-operativo}{%
\subsubsection{Exclusión del panel
operativo}\label{exclusiuxf3n-del-panel-operativo}}

Aunque conceptualmente relevante, la novedad se excluye del panel
principal por fragilidad de validación externa y alta dependencia de
infraestructura (Bornmann, 2019; Fontana et al., 2020).

\hypertarget{mapa-de-cobertura-matemuxe1tica-y-cualificaciuxf3n}{%
\subsection{Mapa de Cobertura Matemática y
Cualificación}\label{mapa-de-cobertura-matemuxe1tica-y-cualificaciuxf3n}}

Para cierre de revisión, explicitamos familias matemáticas y sus límites
de uso:

\begin{longtable}[]{@{}
  >{\raggedright\arraybackslash}p{(\columnwidth - 4\tabcolsep) * \real{0.1379}}
  >{\raggedright\arraybackslash}p{(\columnwidth - 4\tabcolsep) * \real{0.4483}}
  >{\raggedright\arraybackslash}p{(\columnwidth - 4\tabcolsep) * \real{0.4138}}@{}}
\toprule
\begin{minipage}[b]{\linewidth}\raggedright
Familia
\end{minipage} & \begin{minipage}[b]{\linewidth}\raggedright
Referencias principales
\end{minipage} & \begin{minipage}[b]{\linewidth}\raggedright
Cualificación aplicada
\end{minipage} \\
\midrule
\endhead
Rao-Stirling & Porter-Rafols 2009; Rafols-Meyer 2009; Leydesdorff-Rafols
2011 & Sensible a taxonomía y matriz de similitud. \\
Diversidad verdadera & Hill 1973; Jost 2006/2009; Leinster-Cobbold 2012
& Requiere interpretación en escala de números efectivos. \\
Variedad-balance & Shannon/Simpson/Gini; Mutz 2022 & No capturan
disparidad por sí solos. \\
Coherencia \(S\) & Rafols-Meyer 2009; Jensen 2014 & Depende de
decisiones de acoplamiento y umbral. \\
Novedad & Uzzi 2013; Wang 2017; Bornmann 2019; Fontana 2020 & Sensible a
modelo nulo; validación limitada para política. \\
Incertidumbre & Zwanenburg 2022; Nakhoda 2023 & Estimación individual
requiere intervalos de confianza. \\
\bottomrule
\end{longtable}

\hypertarget{muxe1s-alluxe1-de-escalares-enfoques-distribucionales}{%
\subsection{Más allá de escalares: enfoques
distribucionales}\label{muxe1s-alluxe1-de-escalares-enfoques-distribucionales}}

Zhou, Guns y Engels (2023) proponen descomponer flujo interdisciplinario
en triple \((B,I,H)\): \emph{broadness}, \emph{intensity},
\emph{homogeneity}. El resultado es más rico que un escalar único, pero
menos compacto para decisiones operativas.

\hypertarget{resumen-de-taxonomuxeda}{%
\subsection{Resumen de Taxonomía}\label{resumen-de-taxonomuxeda}}

\begin{longtable}[]{@{}lllll@{}}
\toprule
Dimensión & Referencias & Citación & Texto & Red \\
\midrule
\endhead
Diversidad & Rao-Stirling, Shannon, Simpson, Hill & --- & similitud
semántica & --- \\
Coherencia & acoplamiento bibliográfico & --- & --- & centralidad,
clustering \\
Difusión & --- & citas transcampo & --- & --- \\
Novedad & --- & recombinación atípica & parcialmente & modelos nulos \\
\bottomrule
\end{longtable}

Observación central: la celda diversidad/referencias está
sobrerrepresentada frente a coherencia, difusión y validación de
novedad.

\hypertarget{panel-multicomponente}{%
\section{Panel Multicomponente}\label{panel-multicomponente}}

\hypertarget{definiciuxf3n-del-panel}{%
\subsection{Definición del panel}\label{definiciuxf3n-del-panel}}

Trabajamos con el vector:

\[\mathbf P = (\Delta, S, E).\]

No se propone una puntuación compuesta única.

\hypertarget{diversidad-rao-stirling}{%
\subsubsection{Diversidad Rao-Stirling}\label{diversidad-rao-stirling}}

\[\Delta = \sum_{i \neq j}(1-s_{ij})p_ip_j.\]

\hypertarget{coherencia-de-red}{%
\subsubsection{Coherencia de red}\label{coherencia-de-red}}

\[S = \frac{1}{\binom{n}{2}}\sum_{k<l}\cos(\mathbf r_k,\mathbf r_l).\]

\hypertarget{proxy-de-efecto-transcampo}{%
\subsubsection{Proxy de efecto
transcampo}\label{proxy-de-efecto-transcampo}}

\[E = \frac{C_{\text{fuera}}}{C_{\text{total}}}.\]

\hypertarget{discriminaciuxf3n-en-datos-de-juguete}{%
\subsection{Discriminación en datos de
juguete}\label{discriminaciuxf3n-en-datos-de-juguete}}

En el conjunto simulado, integrador y polímata tienen \(\Delta\) similar
pero se separan por \(S\) y \(E\). Especialistas muestran \(\Delta\)
baja y \(S\) alta.

\hypertarget{sensibilidad-y-robustez}{%
\subsection{Sensibilidad y robustez}\label{sensibilidad-y-robustez}}

Para perturbación uniforme de similitud
(\(s_{ij}\to s_{ij}+\varepsilon\)):

\[\Delta_{\text{new}} = \Delta_{\text{old}} - \varepsilon(1-H), \qquad H=\sum_i p_i^2.\]

La discriminación entre perfiles se conserva para perturbaciones
moderadas; la conclusión no depende de un ajuste fino único de
parámetros.

\hypertarget{marco-de-interpretaciuxf3n}{%
\subsection{Marco de interpretación}\label{marco-de-interpretaciuxf3n}}

Patrones típicos:

\begin{itemize}
\tightlist
\item
  \textbf{Integrador}: \(\Delta\) alta, \(S\) alta, \(E\) alta.
\item
  \textbf{Polímata}: \(\Delta\) alta, \(S\) baja, \(E\) baja.
\item
  \textbf{Especialista puente}: \(\Delta\) media, \(S\) alta, \(E\)
  media.
\item
  \textbf{Especialista disciplinar}: \(\Delta\) baja, \(S\) alta, \(E\)
  baja.
\end{itemize}

Las fronteras numéricas deben calibrarse por contexto y reportarse con
incertidumbre.

\hypertarget{mediciuxf3n-en-la-pruxe1ctica}{%
\section{Medición en la Práctica}\label{mediciuxf3n-en-la-pruxe1ctica}}

\hypertarget{autoevaluaciuxf3n-institucional}{%
\subsection{Autoevaluación
institucional}\label{autoevaluaciuxf3n-institucional}}

Las instituciones con CRIS completo pueden computar \(\Delta\) y \(S\)
con datos internos; \(E\) requiere citación externa
(OpenAlex/WoS/Scopus).

\hypertarget{agencias-nacionales-de-evaluaciuxf3n}{%
\subsection{Agencias nacionales de
evaluación}\label{agencias-nacionales-de-evaluaciuxf3n}}

Para categorías como ``Interdisciplinario'', se recomienda composición
de comité guiada por perfil disciplinar observado y no por adscripción
administrativa fija.

\hypertarget{extracciuxf3n-de-datos}{%
\subsection{Extracción de datos}\label{extracciuxf3n-de-datos}}

Requisitos mínimos: referencias estructuradas, clasificación
disciplinar, mapeo de afiliaciones y trazabilidad de citación.

\hypertarget{protocolos-de-mapeo-de-categoruxedas}{%
\subsection{Protocolos de mapeo de
categorías}\label{protocolos-de-mapeo-de-categoruxedas}}

La elección de taxonomía debe documentarse (granularidad, reglas de
multi-asignación, resolución de ambigüedad).

\hypertarget{construcciuxf3n-de-matriz-de-similitud}{%
\subsection{Construcción de matriz de
similitud}\label{construcciuxf3n-de-matriz-de-similitud}}

Opciones: citación histórica, embeddings semánticos o enfoques híbridos.
Cualquier opción exige auditoría de sensibilidad.

\hypertarget{control-de-calidad}{%
\subsection{Control de calidad}\label{control-de-calidad}}

\begin{itemize}
\tightlist
\item
  Cobertura de referencias categorizadas.
\item
  Estabilidad frente a cambios de taxonomía.
\item
  Intervalos de confianza en nivel individual.
\end{itemize}

\hypertarget{implementaciuxf3n-software}{%
\subsection{Implementación software}\label{implementaciuxf3n-software}}

La implementación debe ser reproducible, con separación clara entre
preprocesamiento, cálculo de indicadores y generación de reportes.

\hypertarget{problemas-abiertos-y-luxedneas-futuras}{%
\section{Problemas Abiertos y Líneas
Futuras}\label{problemas-abiertos-y-luxedneas-futuras}}

\begin{enumerate}
\def\labelenumi{\arabic{enumi}.}
\tightlist
\item
  \textbf{Calibración de umbrales por campo y etapa de carrera}.
\item
  \textbf{Modelado de incertidumbre individual} en carteras pequeñas.
\item
  \textbf{Integración de señales semánticas} sin perder
  interpretabilidad.
\item
  \textbf{Medición operativa de transdisciplinariedad} (incluyendo
  impacto no académico y coproducción con actores externos).
\item
  \textbf{Validación externa de indicadores de novedad} para uso
  evaluativo.
\item
  \textbf{Riesgos performativos y manipulación estratégica} bajo
  regímenes de evaluación.
\end{enumerate}

\hypertarget{estudio-de-caso-evaluaciuxf3n-a-nivel-departamental}{%
\section{Estudio de Caso: Evaluación a Nivel
Departamental}\label{estudio-de-caso-evaluaciuxf3n-a-nivel-departamental}}

\hypertarget{contexto-y-datos}{%
\subsection{Contexto y datos}\label{contexto-y-datos}}

Aplicamos el panel a un departamento con siete perfiles simulados,
representando integradores, polímatas, especialistas y casos frontera.

\hypertarget{cuxf3mputo-del-panel}{%
\subsection{Cómputo del panel}\label{cuxf3mputo-del-panel}}

Cada perfil recibe vector (\(\Delta\), \(S\), \(E\)) bajo una taxonomía
común y matriz de similitud auditada.

\hypertarget{interpretaciuxf3n-y-clasificaciuxf3n}{%
\subsection{Interpretación y
clasificación}\label{interpretaciuxf3n-y-clasificaciuxf3n}}

El panel identifica:

\begin{itemize}
\tightlist
\item
  Integradores genuinos (amplitud + integración + difusión).
\item
  Polímatas no integrativos (amplitud sin coherencia).
\item
  Especialistas mal clasificados como interdisciplinarios.
\item
  Casos ambiguos que requieren revisión experta cualitativa.
\end{itemize}

\hypertarget{comparaciuxf3n-con-enfoques-monoindicador}{%
\subsection{Comparación con enfoques
monoindicador}\label{comparaciuxf3n-con-enfoques-monoindicador}}

Métricas únicas (Shannon, Simpson o \(\Delta\) aislada) no discriminan
consistentemente perfiles con implicaciones evaluativas distintas.

\hypertarget{limitaciones}{%
\subsection{Limitaciones}\label{limitaciones}}

Caso ilustrativo, no inferencia poblacional. Las conclusiones deben
validarse con datos institucionales reales y protocolos de calibración
externa.

\hypertarget{conclusiones}{%
\section{Conclusiones}\label{conclusiones}}

La medición bibliométrica de IDR sigue siendo un problema abierto, pero
la revisión muestra una conclusión robusta: ninguna métrica única es
suficiente. El panel multicomponente (\(\Delta\), \(S\), \(E\)) ofrece
una estructura más fiel al fenómeno, mejora capacidad discriminante y
permite decisiones más transparentes.

Este resultado se alinea con las cuatro motivaciones del proyecto:
evaluar impacto/calidad más allá de amplitud superficial, construir
esquemas computables con datos institucionales, distinguir integración
de polimatía y apoyar evaluación justa en agencias nacionales.

La siguiente etapa natural es validación con datos reales, calibración
por campo y formalización de reglas de incertidumbre para despliegue
operativo responsable.

\hypertarget{referencias}{%
\section{Referencias}\label{referencias}}

\begin{itemize}
\tightlist
\item
  Abramo, G., D'Angelo, C. A., and Zhang, L. (2018). A comparison of two
  approaches for measuring interdisciplinary research output: The
  disciplinary diversity of authors vs the disciplinary diversity of the
  reference list. \emph{Journal of Informetrics}, 12(4):1182--1193.
\item
  Aksnes, D. W., Karlstrøm, H., and Piro, F. N. (2026). Self-reported
  and bibliometric interdisciplinarity measures rarely correspond: a
  survey-based comparative analysis of indicators and researcher
  perceptions. \emph{Scientometrics}, 131:189--208.
\item
  Bollen, J., Van de Sompel, H., Hagberg, A., and Chute, R. (2009). A
  principal component analysis of 39 scientific impact measures.
  \emph{PLoS ONE}, 4(6):e6022.
\item
  Borlaug, S. B. and Svartefoss, S. M. (2025). Evaluating
  transdisciplinary research quality. In Sivertsen, G. and Langfeldt, L.
  (eds.), \emph{Challenges in Research Policy}, pp.~13--20. Springer,
  Cham.
\item
  Bornmann, L., Tekles, A., Zhang, H. H., and Ye, F. Y. (2019). Do we
  measure novelty when we analyze unusual combinations of cited
  references? A validation study of bibliometric novelty indicators
  based on F1000Prime data. \emph{Journal of Informetrics},
  13(4):100979.
\item
  Boyack, K. W., Klavans, R., and Börner, K. (2005). Mapping the
  backbone of science. \emph{Scientometrics}, 64(3):351--374.
\item
  Cantone, G. G. (2024). How to measure interdisciplinary research? A
  systemic design for the model of measurement. \emph{Scientometrics},
  129:4937--4982.
\item
  Cantone, G. G. (2025). Estimation of disciplinary similarity with
  large language models. \emph{Scientometrics}, 130(10):5345--5373.
\item
  Choi, B. C. K. and Pak, A. W. P. (2006). Multidisciplinarity,
  interdisciplinarity and transdisciplinarity in health research,
  services, education and policy: 1. Definitions, objectives, and
  evidence of effectiveness. \emph{Clinical and Investigative Medicine},
  29(6):351--364.
\item
  Fontana, M., Iori, M., Montobbio, F., and Sinatra, R. (2020). New and
  atypical combinations: An assessment of novelty and
  interdisciplinarity. \emph{Research Policy}, 49(7):104063.
\item
  Goyanes, M., Demeter, M., Grané, A., Albarrán-Lozano, I., and Gil de
  Zúñiga, H. (2020). A mathematical approach to assess research
  diversity: Operationalization and applicability in communication
  sciences, political science, and beyond. \emph{Scientometrics},
  125(3):2299--2322.
\item
  Hammarfelt, B. (2020). Discipline. \emph{Knowledge Organization},
  47(3):244--256.
\item
  Hill, M. O. (1973). Diversity and evenness: A unifying notation and
  its consequences. \emph{Ecology}, 54(2):427--432.
\item
  Jensen, P. and Lutkouskaya, K. (2014). The many dimensions of
  laboratories' interdisciplinarity. \emph{Scientometrics},
  98(1):619--631.
\item
  Jost, L. (2006). Entropy and diversity. \emph{Oikos}, 113(2):363--375.
\item
  Jost, L. (2009). Mismeasuring biological diversity: Response to
  Hoffmann and Hoffmann (2008). \emph{Ecological Economics},
  68(4):925--928.
\item
  Klein, J. T. (2008). Evaluation of interdisciplinary and
  transdisciplinary research: A literature review. \emph{American
  Journal of Preventive Medicine}, 35(2S):S116--S123.
\item
  Larivière, V. and Gingras, Y. (2010). On the relationship between
  interdisciplinarity and scientific impact. \emph{Journal of the
  American Society for Information Science and Technology},
  61(1):126--131.
\item
  Leinster, T. and Cobbold, C. A. (2012). Measuring diversity: the
  importance of species similarity. \emph{Ecology}, 93(3):477--489.
\item
  Leydesdorff, L. (2006). Can scientific journals be classified in terms
  of aggregated journal--journal citation relations using the Journal
  Citation Reports? \emph{Journal of the American Society for
  Information Science and Technology}, 57(5):601--613.
\item
  Leydesdorff, L. and Rafols, I. (2011). Indicators of the
  interdisciplinarity of journals: Diversity, centrality, and citations.
  \emph{Journal of Informetrics}, 5(1):87--100.
\item
  Leydesdorff, L., Wagner, C. S., and Bornmann, L. (2019).
  Interdisciplinarity as diversity in citation patterns among journals:
  Rao-Stirling diversity, relative variety, and the Gini coefficient.
  \emph{Journal of Informetrics}, 13(1):255--269.
\item
  Marres, N. and de Rijcke, S. (2020). From indicators to indicating
  interdisciplinarity: A participatory mapping methodology for research
  communities in-the-making. \emph{Quantitative Science Studies},
  1(3):1041--1055.
\item
  Morillo, F., Bordons, M., and Gómez, I. (2001). An approach to
  interdisciplinarity through bibliometric indicators.
  \emph{Scientometrics}, 51(1):203--222.
\item
  Morillo, F., Bordons, M., and Gómez, I. (2003). Interdisciplinarity in
  science: A tentative typology of disciplines and research areas.
  \emph{Journal of the American Society for Information Science and
  Technology}, 54(13):1237--1249.
\item
  Moulton, R. and Jiang, Y. (2018). Maximally consistent sampling and
  the Jaccard index of probability distributions. In \emph{Proceedings
  of the 2018 IEEE International Conference on Data Mining (ICDM)},
  pages 347--356. IEEE.
\item
  Mutz, R. (2022). Diversity and interdisciplinarity: Should variety,
  balance and disparity be combined as a product or better as a sum? An
  information-theoretical and statistical estimation approach.
  \emph{Scientometrics}, 127(12):7397--7414.
\item
  Nakhoda, M., Whigham, P., and Zwanenburg, S. (2023). Quantifying and
  addressing uncertainty in the measurement of interdisciplinarity.
  \emph{Scientometrics}, 128:6107--6127.
\item
  Noji, H., Yasuda, R., Yoshida, M., and Kinosita, K. (1997). Direct
  observation of the rotation of F1-ATPase. \emph{Nature},
  386(6622):299--302.
\item
  OECD (1998). \emph{Interdisciplinarity in Science and Technology}.
  Paris: OECD Directorate for Science, Technology and Industry.
\item
  Porter, A. L. and Rafols, I. (2009). Is science becoming more
  interdisciplinary? Measuring and mapping six research fields over
  time. \emph{Scientometrics}, 81(3):719--745.
\item
  Rafols, I. (2019). S\&T indicators in the wild: contextualization and
  participation for responsible metrics. \emph{Research Evaluation},
  28(1):7--22.
\item
  Rafols, I. and Meyer, M. (2009). Diversity and network coherence as
  indicators of interdisciplinarity: case studies in bionanoscience.
  \emph{Scientometrics}, 82(2):263--287.
\item
  Stirling, A. (2007). A general framework for analysing diversity in
  science, technology and society. \emph{Journal of the Royal Society
  Interface}, 4(15):707--719.
\item
  Stokols, D., Fuqua, J., Gress, J., Harvey, R., Phillips, K.,
  Baezconde-Garbanati, L., Unger, J., Palmer, P., Clark, M. A., Colby,
  S. M., Morgan, G., and Trochim, W. (2003). Evaluating
  transdisciplinary science. \emph{Nicotine \& Tobacco Research},
  5(Suppl 1):S21--S39.
\item
  Tomishige, M., Klopfenstein, D. R., and Vale, R. D. (2002). Conversion
  of Unc104/KIF1A kinesin into a processive motor after dimerization.
  \emph{Science}, 297(5590):2263--2267.
\item
  Uzzi, B., Mukherjee, S., Stringer, M., and Jones, B. (2013). Atypical
  combinations and scientific impact. \emph{Science},
  342(6157):468--472.
\item
  Wang, J., Veugelers, R., and Stephan, P. (2017). Bias against novelty
  in science: A cautionary tale for users of bibliometric indicators.
  \emph{Research Policy}, 46(8):1416--1436.
\item
  Wagner, C. S., Roessner, J. D., Bobb, K., Klein, J. T., Boyack, K. W.,
  Keyton, J., Rafols, I., and Börner, K. (2011). Approaches to
  understanding and measuring interdisciplinary scientific research
  (IDR): A review of the literature. \emph{Journal of Informetrics},
  5(1):14--26.
\item
  Wang, Q. and Schneider, J. W. (2020). Consistency and validity of
  interdisciplinarity measures. \emph{Quantitative Science Studies},
  1(1):239--263.
\item
  Xiang, S., Romero, D. M., and Teplitskiy, M. (2025). Evaluating
  interdisciplinary research: Disparate outcomes for topic and knowledge
  base. \emph{Proceedings of the National Academy of Sciences},
  122(16):e2409752122.
\item
  Zhang, L., Rousseau, R., and Glänzel, W. (2016). Diversity of
  references as an indicator of the interdisciplinarity of journals:
  Taking similarity between subject fields into account. \emph{Journal
  of the Association for Information Science and Technology},
  67(5):1257--1265.
\item
  Zhou, Q., Guns, R., and Engels, T. C. E. (2023). Towards indicating
  interdisciplinarity: Characterizing interdisciplinary knowledge flow.
  \emph{Journal of the Association for Information Science and
  Technology}, 74(11):1325--1340.
\item
  Zwanenburg, S., Nakhoda, M., and Whigham, P. (2022). Toward greater
  consistency and validity in measuring interdisciplinarity: a
  systematic and conceptual evaluation. \emph{Scientometrics},
  127:3035--3065.
\end{itemize}

\end{document}
